\subsection{显示偏好关系的闭包证明}

\begin{proposition}[传递闭包的最小性]
设 $R$ 是有限方案上的直接显示偏好关系，记 $R^+$ 为所有有限有向路径生成的关系。则 $R^+$ 是包含 $R$ 的最小传递关系。
\end{proposition}
\begin{proof}
任意 $(x,y)\in R$ 是长度为一的路径，故 $R\subseteq R^+$。若 $(x,y),(y,z)\in R^+$，分别存在从 $x$ 到 $y$、从 $y$ 到 $z$ 的有限路径，拼接两条路径即得从 $x$ 到 $z$ 的有限路径，因此 $R^+$ 传递。若 $S$ 是包含 $R$ 的任意传递关系，则对路径长度作归纳：长度一的边在 $S$ 中；若长度 $k$ 的路径对应关系在 $S$ 中，再接一条 $R$ 边，由传递性仍在 $S$ 中。故 $R^+\subseteq S$。
\end{proof}

\begin{example}[闭包如何发现隐含矛盾]
若直接关系只有 $a\succeq^Rb$、$b\succeq^Rc$，则闭包增加 $a\succeq^Rc$。若另有严格边 $c\succ a$，闭包便产生严格环。逐项查看原始边无法发现该矛盾，必须先计算闭包。
\end{example}

\subsection{HARP 检验中的价格—支出不等式}
对观测 $t,s$，若 $p^t\cdot x^s\le p^t\cdot x^t$，则 $x^s$ 在 $t$ 期可负担。若同时 $p^s\cdot x^t<p^s\cdot x^s$，则 $x^t$ 在 $s$ 期严格不可负担，形成严格反向显示偏好。HARP 正是要求这类价格—支出不等式不能沿路径闭合。

\begin{proposition}[支出不等式的尺度不变性]
若每期价格与收入同时乘以同一正数 $a_t$，则 GARP/HARP 的可负担关系不变。
\end{proposition}
\begin{proof}
$a_tp^t\cdot x^s\le a_tp^t\cdot x^t$ 当且仅当 $p^t\cdot x^s\le p^t\cdot x^t$，因为 $a_t>0$ 不改变不等式方向。故所有直接显示偏好边及其闭包均不变。
\end{proof}

\begin{remark}
该尺度不变性解释了为什么显示偏好检验只依赖相对价格与相对支出，而不依赖货币单位。把元改为分、把收入和价格同时乘以一百，不会改变理性化结论。
\end{remark}
